% =============================================================================
% kovex_paper.tex  —  FINAL VERSION with real agent effectiveness experiment
% Kovex: Graph-Aware, Per-Node Coordination for Multi-Agent LLM Coding Systems
% =============================================================================

\documentclass[12pt]{article}

\usepackage[margin=1.25in]{geometry}
\usepackage[utf8]{inputenc}
\usepackage[T1]{fontenc}
\usepackage{amsmath, amssymb, amsthm}
\usepackage{booktabs}
\usepackage{array}
\usepackage{multirow}
\usepackage{hyperref}
\usepackage[numbers, sort&compress]{natbib}
\usepackage{setspace}
\usepackage{enumitem}
\usepackage{mdframed}
\usepackage{tcolorbox}
\tcbuselibrary{skins, breakable}
\usepackage{algorithm}
\usepackage{algpseudocode}
\usepackage{tikz}
\usetikzlibrary{arrows.meta, positioning, shapes.geometric, fit, backgrounds, calc}
\usepackage{xcolor}
\usepackage{siunitx}

\hypersetup{colorlinks=true, citecolor=blue!60!black,
            urlcolor=blue!60!black, linkcolor=blue!60!black}

\newtheorem{definition}{Definition}
\tcbset{defbox/.style={colback=blue!4!white, colframe=blue!40!black,
  fonttitle=\bfseries, sharp corners=south, boxrule=0.6pt,
  leftrule=3pt, before skip=1em, after skip=1em}}
\newtcolorbox[auto counter]{kovexdef}[1][]{defbox,
  title={Definition~\thetcbcounter}, #1}

\algnewcommand\algorithmicforeach{\textbf{for each}}
\algdef{S}[FOR]{ForEach}[1]{\algorithmicforeach\ #1\ \textbf{do}}

\begin{document}
\onehalfspacing

\begin{center}
  {\LARGE\bfseries Kovex: Graph-Aware, Per-Node Coordination\\[4pt]
  for Multi-Agent LLM Coding Systems}\\[1.4em]
  {\large Ibrahem Alkurdi}\\[0.3em]
  {\normalsize \texttt{ibrahem.alkurdi@gmail.com}}\\[1.0em]
  {\normalsize \today}
\end{center}

\begin{abstract}
Multi-agent LLM coding systems operating concurrently on shared codebases face
a systematic failure mode that neither message-passing coordination nor
text-level concurrency control can address: \emph{semantic conflicts}, arising
when one agent's write to a function signature or interface boundary invalidates
assumptions held by a dependent agent. We present \textbf{Kovex}, a coordination
substrate that indexes a live codebase into a Neo4j property graph, exposes it
to agents via a Model Context Protocol (MCP) server, and notifies registered
dependents via bounded breadth-first search over the call graph when a write
is committed. We evaluate Kovex in two stages. A systems correctness evaluation
on 14 tasks confirms notifications fire at the correct call-graph distances
($d \in \{0,1,2\}$) with 100\% accuracy and mean latency 45.6\,ms, while
correctly suppressing notifications at $d=3$ ($k=2$ false-negative rate 14.3\%).
A real agent effectiveness experiment using Claude Sonnet~4.6 as both agents
(50 trials, $n=5$ per cell) shows that Kovex-coordinated agents achieve an
80\% final-code pass rate on tasks within the notification radius versus 0\%
for uncoordinated baseline agents, with two tasks reaching 100\% vs.\ 0\%.
The T13 false-negative control correctly collapses to near-baseline (20\% vs.\ 0\%,
unattributable to Kovex), validating the CALLS-only BFS design.
\end{abstract}

% =============================================================================
\section{Introduction}
% =============================================================================

The deployment of large language models as collaborative coding agents has
advanced rapidly, yet a fundamental coordination problem has emerged: when two
agents operate concurrently on a shared codebase, one agent's modification to a
function signature or type definition can silently invalidate the code being
generated by the other. CooperBench~\cite{khatua2026cooperbench} quantifies this
as the \emph{curse of coordination}, reporting a 30\% lower success rate for
collaborative agent pairs compared to single-agent execution. CodeCRDT~\cite{pugachev2025codecrdt}
demonstrates that text-level concurrency control, while guaranteeing syntactic
convergence, still produces a 5--10\% semantic conflict rate. Existing frameworks
resolve concurrency by eliminating it (MetaGPT~\cite{hong2024metagpt},
ChatDev~\cite{qian2024chatdev}) or by deferring to global test re-execution
(AgentForge~\cite{kumar2026agentforge}).

We identify the root cause as a missing representation layer. No system maintains
a persistent, queryable model of the codebase's dependency structure through which
agents can register read dependencies and receive targeted notifications when those
dependencies change. Without this layer, verification is either absent or exhaustive.

We contribute \textbf{Kovex}, comprising: (1) a Neo4j property graph indexed from
live source code; (2) an MCP server through which agents register read dependencies
on named graph nodes; and (3) a write-triggered bounded BFS notification algorithm
that delivers targeted re-verification directives to registered dependents within
a configurable $k$-hop call-graph radius. Our central research question is:

\begin{center}
\emph{Does graph-aware, per-node read/write coordination with write-triggered}\\
\emph{targeted verification reduce semantic conflict rates in multi-agent coding}\\
\emph{systems, relative to uncoordinated baselines?}
\end{center}

We answer this with two complementary evaluations: a systems correctness study
confirming the coordination mechanism operates as specified, and a real agent
effectiveness experiment with LLM agents demonstrating a measurable downstream
reduction in semantic conflict rates.

% =============================================================================
\section{Related Work}
% =============================================================================

\subsection{Multi-Agent LLM Systems for Software Engineering}

MetaGPT~\cite{hong2024metagpt} encodes Standardized Operating Procedures into
prompt sequences, assigning specialized roles that hand off work through structured
artifacts. ChatDev~\cite{qian2024chatdev} implements a similar role hierarchy
through natural-language dialogue. AutoGen~\cite{wu2023autogen} generalizes this
into a flexible multi-agent conversation framework. All three enforce coordination
sequentially; none maintains a runtime dependency model of the codebase.
SWE-agent~\cite{yang2024sweagent} focuses on single-agent performance on
SWE-bench~\cite{jimenez2024swebench}; SWE-EVO~\cite{thai2025sweevo} extends
evaluation to long-horizon tasks requiring coordinated modifications across an
average of 21 files, where the strongest agents achieve only 21\% resolution.

CooperBench~\cite{khatua2026cooperbench} provides the most direct empirical
characterization of the coordination problem. Across 652 collaborative coding
tasks, agent pairs achieve 30\% lower success rates than single agents, with
communication reducing syntactic merge conflicts but failing to prevent semantic
conflicts arising from interface incompatibilities.
AgentForge~\cite{kumar2026agentforge} advances the state of the art by binding
verification to each state transition via sandboxed execution, achieving 40.0\%
on SWE-bench Lite. However, its verification is global and cannot determine
from the codebase's dependency structure which agents need to re-verify.

\subsection{Graph-Based Code Intelligence}

Ferrante et al.~\cite{ferrante1987pdg} introduced the Program Dependence Graph,
unifying control and data dependencies for each program operation.
Yamaguchi et al.~\cite{yamaguchi2014cpg} unified the AST, CFG, and PDG into the
Code Property Graph stored in a graph database, demonstrating that code structure
can serve as a first-class database artifact rather than a transient in-memory
representation. ACER~\cite{chen2023acer} provides a language-agnostic pipeline
using Tree-sitter~\cite{utkin2022parsers} to extract call graphs at scale.
GraphRAG~\cite{edge2024graphrag} and its multi-agent
extension~\cite{gusarov2025multiagentgraphrag} show that LLMs can reason
effectively over graph-structured knowledge. None of these uses a code graph
as a live write-aware coordination layer.

\subsection{Concurrency Control and Collaborative Editing}

Ellis and Gibbs~\cite{ellis1989ot} introduced Operational Transformation for
convergent collaborative text editing. Shapiro et al.~\cite{shapiro2011crdts}
formalized CRDTs, guaranteeing strong eventual consistency without locking.
Pugachev~\cite{pugachev2025codecrdt} applies CRDTs to multi-agent LLM code
generation, achieving 100\% syntactic convergence across 600 trials while
reporting a residual 5--10\% semantic conflict rate. This establishes that
text-level consistency is strictly weaker than the coordination guarantee
Kovex targets. Gray et al.~\cite{gray1975locks} analyze lock granularity
trade-offs; Kovex adopts an analogous design at semantic node granularity.

\subsection{Verification, Semantic Conflicts, and Interface Contracts}

Meyer's Design by Contract~\cite{meyer1992dbc} established function signatures
as contractual obligations between implementors and callers. Krieter et
al.~\cite{krieter2020mergeconflicts} show that higher-order semantic merge
conflicts are common in open-source projects even without concurrent agents.
Santos et al.~\cite{santos2024semanticconflicts} and Barbosa et
al.~\cite{barbosa2025pointeranalysis} propose post-hoc static analysis for
detection. Kovex targets the same failure class but prevents conflicts at
coordination time. Consumer-driven contract testing~\cite{lehva2019contracts,
aalto2024contracts} provides a structural analogy: a read-registration is
equivalent to a consumer contract, and write-triggered notification is the
automated provider responsibility mechanism.

% =============================================================================
\section{System Design}
% =============================================================================

\subsection{Overview}

Kovex comprises three components: the \emph{Kovex Graph} (a Neo4j labeled
property graph indexed from live source code), the \emph{MCP Server} (a
Model Context Protocol server exposing the graph to agents as a tool API and
maintaining a read registry), and the \emph{Write Log} (an append-only record
of every write event and its notifications).

\subsection{Graph Schema}

The graph models the \emph{semantic contract surface} of a codebase. Four
node labels are defined: \texttt{Function} (properties: \texttt{name},
\texttt{file}, \texttt{signature}, \texttt{return\_type}, \texttt{body\_hash}),
\texttt{Type}, \texttt{Module}, and \texttt{Symbol}. Seven edge types encode
semantic dependencies: \textsc{Calls}, \textsc{Uses\_Type}, \textsc{Defines},
\textsc{Imports}, \textsc{Inherits}, \textsc{Implements}, and \textsc{Overrides}.
The graph is indexed using Tree-sitter~\cite{utkin2022parsers} for
language-agnostic AST extraction.

\subsection{Read/Write State Model}

Let $\mathcal{G} = (\mathcal{V}, \mathcal{E})$ be the Kovex graph and
$\mathcal{A}$ the set of active agents. Each agent $a$ maintains a read set
$\mathrm{RS}(a) \subseteq \mathcal{V}$. The global read registry is
$\mathcal{R} = \{(a, v, t_r, h_r) \mid v \in \mathrm{RS}(a)\}$. A write by
agent $a$ to node $v$ is \emph{observable} iff $h_{\mathrm{post}}(v) \neq
h_{\mathrm{pre}}(v)$; non-observable writes do not trigger notifications.

\subsection{Write-Triggered Notification Algorithm}

\begin{kovexdef}[title={Definition~1: $k$-Hop Notification Radius}]
Let $d_{\mathcal{G}}(u,v)$ be the shortest-path distance restricted to
\textsc{Calls} edges. The \emph{$k$-hop dependency neighborhood} of $v$ is
$\mathcal{N}_k(v) = \{u \in \mathcal{V} \mid d_{\mathcal{G}}(v,u) \leq k\}$.
Agent $a$ is a notification target for a write to $v$ iff
$\mathrm{RS}(a) \cap \mathcal{N}_k(v) \neq \emptyset$.
\end{kovexdef}

\noindent\textbf{Edge type selection.}
Restricting BFS to \textsc{Calls} edges ensures notifications fire strictly on
caller/callee relationships. An all-edge-type BFS places every co-module
function pair at $d \leq 2$ via \textsc{Defines} edges, conflating module
co-location with semantic dependency. Algorithm~\ref{alg:notify} gives the
formal procedure.

\begin{algorithm}[ht]
\caption{Write-Triggered Notification (\textsc{NotifyDependents})}
\label{alg:notify}
\small
\begin{algorithmic}[1]
\Require Graph $\mathcal{G}$, written node $v$, radius $k$, registry $\mathcal{R}$
\Ensure Notification triples $(a, u, d)$
\State $\textit{frontier} \leftarrow \{v\}$;\;
       $\textit{visited} \leftarrow \{v\}$;\;
       $\textit{notifs} \leftarrow \emptyset$
\For{$d = 1$ \textbf{to} $k$}
    \State $\textit{next} \leftarrow \emptyset$
    \ForEach{$u \in \textit{frontier}$}
        \ForEach{$w \in \mathrm{CallNeighbors}(u,\mathcal{G})$,\; $w \notin \textit{visited}$}
            \State $\textit{visited} \mathrel{+}= w$;\; $\textit{next} \mathrel{+}= w$
            \ForEach{$(a,w,t_r,h_r) \in \mathcal{R}$}
                \State $\textit{notifs} \mathrel{+}= (a,w,d)$
            \EndFor
        \EndFor
    \EndFor
    \State $\textit{frontier} \leftarrow \textit{next}$
\EndFor
\State \Return $\textit{notifs}$
\end{algorithmic}
\end{algorithm}

\subsection{MCP Server Interface}

Eight tools are exposed: \texttt{register\_read}, \texttt{deregister\_read},
\texttt{list\_reads}, \texttt{commit\_write} (applies patch, runs
Algorithm~\ref{alg:notify}, logs event, returns notify set),
\texttt{get\_node} (ephemeral read; no registry side effect),
\texttt{get\_dependents}, \texttt{query\_graph} (read-only Cypher), and
\texttt{poll\_directives}. The distinction between \texttt{get\_node} and
\texttt{register\_read} is a deliberate interface contract: ephemeral reads
do not enroll the agent in the notification protocol.

\subsection{Semantic Conflict Definition}

\begin{kovexdef}[title={Definition~2: Semantic Conflict}]
Let $\mathrm{Valid}(C, \mathcal{G}_t)$ denote that code $C$ type-checks against
codebase state $\mathcal{G}_t$. A \textbf{semantic conflict} between agents $A$
and $B$ with respect to node $v$ at time $t_w$ holds iff:
\begin{enumerate}[label=\textnormal{(SC\arabic*)}, leftmargin=3em, itemsep=0pt]
  \item Agent $A$ commits an observable write to $v$ at $t_w$.
  \item $(B,u,t_r,h_r) \in \mathcal{R}$ for some $u \in \mathcal{N}_k(v)$,
        $t_r < t_w$.
  \item Agent $B$ produced code $C_B$ in $[t_r, t_w)$.
  \item $\mathrm{Valid}(C_B,\mathcal{G}_{t_w^-}) \wedge
        \neg\,\mathrm{Valid}(C_B,\mathcal{G}_{t_w^+})$.
\end{enumerate}
\end{kovexdef}

% =============================================================================
\section{Evaluation}
% =============================================================================

\subsection{Codebase and Graph Statistics}

We index the httpx Python HTTP client library~\cite{httpx}. The resulting
Kovex graph contains 550 nodes (358 \texttt{Function}, 85 \texttt{Type},
84 \texttt{Symbol}, 23 \texttt{Module}) and 916 edges (527 \textsc{Defines},
225 \textsc{Uses\_Type}, 99 \textsc{Calls}, 56 \textsc{Inherits},
9 \textsc{Imports}).

\subsection{Stage 1: Systems Correctness Evaluation}

We design 14 evaluation tasks organized by CALLS-chain distance $d$ between
write target and Agent~B's registration node, covering three distance groups
and a false-negative control group. Each task applies a realistic signature
change (parameter removal, renaming, or addition) and verifies Definition~2
condition~(SC4) via \texttt{mypy}. Table~\ref{tab:tasks} summarizes all tasks.

\begin{table}[ht]
\centering
\caption{Evaluation tasks. $d$: CALLS-chain distance. FN: false-negative control.}
\label{tab:tasks}
\small
\renewcommand{\arraystretch}{1.2}
\begin{tabular}{@{}llllc@{}}
\toprule
\textbf{Task} & \textbf{Write Target} & \textbf{B Registration}
              & \textbf{Change} & \textbf{$d$} \\
\midrule
T01 & \texttt{AsyncClient.send}         & \texttt{AsyncClient.send}         & param removal  & 0 \\
T02 & \texttt{AsyncClient.request}      & \texttt{AsyncClient.request}      & param rename   & 0 \\
\midrule
T03 & \texttt{AsyncClient.send}         & \texttt{AsyncClient.request}      & param removal  & 1 \\
T04 & \texttt{AsyncClient.send}         & \texttt{AsyncClient.\_send\_handling\_auth} & param removal & 1 \\
T05 & \texttt{AsyncClient.request}      & \texttt{AsyncClient.get}          & param rename   & 1 \\
T06 & \texttt{AsyncClient.request}      & \texttt{AsyncClient.post}         & param rename   & 1 \\
T07 & \texttt{AsyncClient.\_send\_handling\_redirects} & \texttt{AsyncClient.\_send\_handling\_auth} & param addition & 1 \\
T08 & \texttt{AsyncClient.\_send\_single\_request} & \texttt{AsyncClient.\_send\_handling\_redirects} & param addition & 1 \\
\midrule
T09 & \texttt{AsyncClient.send}         & \texttt{AsyncClient.get}          & param removal  & 2 \\
T10 & \texttt{AsyncClient.send}         & \texttt{AsyncClient.post}         & param removal  & 2 \\
T11 & \texttt{AsyncClient.\_send\_handling\_redirects} & \texttt{AsyncClient.send} & param addition & 2 \\
T12 & \texttt{AsyncClient.\_send\_single\_request} & \texttt{AsyncClient.\_send\_handling\_auth} & param addition & 2 \\
\midrule
T13 & \texttt{AsyncClient.\_send\_single\_request} & \texttt{AsyncClient.send}   & param addition & 3 (FN) \\
T14 & \texttt{AsyncClient.\_transport\_for\_url}   & \texttt{AsyncClient.\_send\_handling\_auth} & param addition & 3 (FN) \\
\bottomrule
\end{tabular}
\end{table}

All 14 tasks pass every applicable check: 12/12 tasks at $d \leq 2$ receive
exactly one directive with the correct distance label; T13 and T14 ($d=3 > k=2$)
receive zero directives. All 14 baseline runs receive zero directives.
Table~\ref{tab:latency} reports notification latency by distance group.

\begin{table}[ht]
\centering
\caption{Notification latency by distance group ($k=2$, CALLS-only BFS).}
\label{tab:latency}
\small
\renewcommand{\arraystretch}{1.2}
\begin{tabular}{@{}lcccc@{}}
\toprule
\textbf{Group} & \textbf{$n$} & \textbf{Mean (ms)} & \textbf{Range (ms)} & \textbf{Std (ms)} \\
\midrule
$d = 0$ & 2  & 39.2 & [34.0,\;44.3] & 7.3  \\
$d = 1$ & 6  & 47.3 & [27.9,\;81.8] & 21.2 \\
$d = 2$ & 4  & 46.2 & [32.9,\;65.5] & 14.7 \\
\midrule
All     & 12 & 45.6 & [27.9,\;81.8] & 17.4 \\
\bottomrule
\end{tabular}
\end{table}

All 12 notifications arrive within 500\,ms (maximum 81.8\,ms). Mean latency
does not increase monotonically with $d$, reflecting that Neo4j query overhead
dominates BFS depth at this codebase scale.

\subsection{Stage 2: Real Agent Effectiveness Experiment}

\paragraph{Setup.}
To measure whether Kovex notifications causally improve agent code quality,
we run a real two-agent experiment using the Anthropic Claude Sonnet~4.6 model
as both Agent~A (writer) and Agent~B (reader). We select five representative
tasks: T01 ($d=0$), T03 ($d=1$), T05 ($d=1$), T09 ($d=2$), and T13 ($d=3$,
false-negative control). Each task is run $n=5$ times per condition
(50 trials total).

\paragraph{Protocol.}
In the \emph{Kovex condition}, Agent~B's Phase~1 session calls
\texttt{register\_read} on its registration node, queries the current signature
via \texttt{get\_node}, and generates initial code. Agent~A then calls
\texttt{commit\_write}. Agent~B's Phase~2 session---initialized with Phase~1
context---calls \texttt{poll\_directives} to receive any Kovex notification,
calls \texttt{get\_node} to read the updated signature, and revises its code.
In the \emph{baseline condition}, Agent~B does not call \texttt{register\_read}
or \texttt{poll\_directives} and submits its Phase~1 code unchanged after
Agent~A's write. The metric is \texttt{mypy} pass rate on Agent~B's final code
against the post-write codebase, filtered to errors in Agent~B's file only
(\texttt{agent\_code\_post\_passed}).

\paragraph{Results.}
Table~\ref{tab:agent} reports per-task pass rates across five runs.

\begin{table}[ht]
\centering
\caption{Real agent effectiveness results ($n=5$ per cell, Claude Sonnet~4.6).
Kovex (notified) vs.\ Baseline (uncoordinated). T13 is the false-negative
control: Kovex correctly issues no notification ($d=3 > k=2$).}
\label{tab:agent}
\small
\renewcommand{\arraystretch}{1.3}
\begin{tabular}{@{}llccc@{}}
\toprule
\textbf{Task} & \textbf{Group} & \textbf{Kovex Pass} & \textbf{Baseline Pass} & \textbf{$\Delta$} \\
\midrule
T01 & $d=0$ & 4/5 (80\%)  & 0/5 (0\%)  & +80pp \\
T03 & $d=1$ & 5/5 (100\%) & 0/5 (0\%)  & +100pp \\
T05 & $d=1$ & 4/5 (80\%)  & 0/5 (0\%)  & +80pp \\
T09 & $d=2$ & 3/5 (60\%)  & 0/5 (0\%)  & +60pp \\
\midrule
T01--T09 aggregate & $d \leq 2$ & 16/20 (80\%) & 0/20 (0\%) & +80pp \\
\midrule
T13 & $d=3$ (FN) & 1/5 (20\%)$^*$ & 0/5 (0\%) & $\approx$0pp \\
\bottomrule
\multicolumn{5}{@{}l@{}}{$^*$ All 5 T13 Kovex trials confirmed \texttt{notify\_set=[]} ---}\\
\multicolumn{5}{@{}l@{}}{the pass is unattributable to Kovex (agent wrote compatible code by chance).}\\
\end{tabular}
\end{table}

On the 20 trials where Kovex is expected to notify ($d \leq 2$),
coordinated agents pass at 80\% vs.\ 0\% for uncoordinated baseline agents
--- a gap of 80 percentage points. T03 and T05 achieve a perfect 5/5 vs.\ 0/5
split. T09 ($d=2$) achieves 60\% vs.\ 0\%, demonstrating that the two-hop
traversal produces actionable directives even at transitive distance.

The T13 false-negative control validates the CALLS-only BFS design: Kovex
issues zero notifications across all 5 runs ($d=3 > k=2$), and the Kovex
pass rate of 20\% is within noise of the 0\% baseline, confirming that
Agent~B does not benefit from coordination when Kovex correctly withholds
the notification.

\paragraph{Failure analysis.}
The remaining 4/20 Kovex failures are attributable to CLI tool-protocol
brittleness: the model occasionally emits a \texttt{<tool\_call>} block and
a \texttt{<python\_code>} block in the same turn; truncation discards the
code block, and the subsequent turn treats the (now absent) code as final.
This failure mode is independent of Kovex's notification mechanism and would
be eliminated by native SDK tool use, which enforces strict turn structure.
The 80\% figure is therefore a conservative lower bound on Kovex's
effectiveness under a reliable tool-use protocol.

% =============================================================================
\section{Discussion}
% =============================================================================

\paragraph{Edge type selection as a design parameter.}
An initial implementation used all edge types (\textsc{Calls},
\textsc{Uses\_Type}, \textsc{Defines}), placing every co-module function pair
at $d \leq 2$ via \textsc{Defines} edges. This eliminates false negatives for
same-module writes but conflates module co-location with semantic dependency.
Restricting BFS to \textsc{Calls} edges restores intended semantics:
notifications fire strictly on caller/callee relationships. This design choice
should be made explicit in any deployment.

\paragraph{Sensitivity to $k$.}
The 14.3\% false-negative rate at $k=2$ (T13, T14) would be eliminated by
increasing $k$ to~3 at the cost of additional false-positive notifications.
The agent experiment confirms this trade-off is meaningful: T13's Kovex pass
rate of 20\% matches the baseline, showing that withheld notifications produce
no benefit. Optimal $k$ is a function of codebase topology and the relative
cost of false positives versus false negatives.

\paragraph{Comparison to related work.}
Kovex addresses the 5--10\% residual semantic conflict rate of
CodeCRDT~\cite{pugachev2025codecrdt} by operating at the semantic graph layer.
It addresses AgentForge's~\cite{kumar2026agentforge} global verification
overhead by restricting re-verification to the $k$-hop call neighborhood.
The two approaches are complementary: Kovex issues targeted directives, and
AgentForge-style sandboxed execution can serve as the verification mechanism
triggered by those directives.

\paragraph{Limitations.}
The experiment uses a two-agent, single-codebase setting with deterministic
signature changes and a CLI tool protocol that introduces noise. Real
deployments involve more concurrent agents, non-deterministic generation, and
simultaneous multi-node writes. The conflict classification mechanism is
Python-specific. Scaling to production multi-agent systems and additional
languages remains future work.

% =============================================================================
\section{Conclusion and Future Work}
% =============================================================================

We presented Kovex, a graph-aware coordination substrate for multi-agent LLM
coding systems. The core insight is that a codebase's call-graph dependency
structure, represented as a queryable property graph and exposed via an MCP
server, can serve as a coordination layer: agents register read dependencies
on named nodes, and writes trigger targeted verification directives to all
registered dependents within a $k$-hop call-graph radius.

Systems correctness evaluation on 14 tasks demonstrates 100\% notification
accuracy for $d \leq 2$ (mean latency 45.6\,ms) and a 14.3\% false-negative
rate at $d=3$ for $k=2$. Real agent effectiveness evaluation with
Claude~Sonnet~4.6 shows coordinated agents achieve 80\% final-code pass rate
versus 0\% for uncoordinated agents on tasks within the notification radius,
with two tasks reaching 100\% vs.\ 0\%. The false-negative control correctly
collapses to near-baseline, validating the CALLS-only BFS design.

We identify six directions for future work.

\textbf{Native SDK tool use.}
The current agent experiment uses a CLI-based tool protocol that introduces
noise from hallucinated tool results, capping the measured effectiveness at
80\% as a conservative lower bound. Re-running the same 50 trials with the
Anthropic Python SDK, which enforces strict turn structure and prevents
hallucinated tool results, would either confirm the 80\% figure or reveal
a higher true upper bound --- both outcomes are informative.

\textbf{Optimal $k$-hop study.}
The choice of $k=2$ is currently fixed without empirical justification.
A systematic study running the agent experiment at $k \in \{1,2,3,4\}$
across all five tasks would yield a three-way trade-off curve: as $k$
increases, pass rate and recall rise while notification volume and false-positive
rate grow. Fitting a simple model to this curve, ideally across codebases of
varying call-graph density, would produce an empirically grounded recommendation
for $k$ as a function of codebase topology --- a deployable guideline absent
from the current work.

\textbf{Notification precision measurement.}
The current evaluation measures recall (did Kovex notify when it should?) but
not precision (did every notification correspond to a real conflict?). All
current tasks are designed so every notification is a true positive by
construction. Adding tasks where Agent~A makes a semantically safe change
--- such as adding an optional parameter with a default value --- and measuring
how often Kovex notifies Agent~B unnecessarily would complete the precision/recall
picture and establish a more rigorous characterization of the false-positive rate.

\textbf{Cross-codebase generalizability.}
All results are obtained on a single codebase (httpx). A reviewer will
reasonably ask whether the findings generalize to codebases with different
structural properties --- denser call graphs, deeper inheritance hierarchies,
or more cross-module dependencies. Replicating the agent experiment on a
structurally distinct codebase (e.g., \texttt{requests} for a more centralized
call graph, or \texttt{fastapi} for a richer type hierarchy) would establish
whether the effectiveness gap holds across codebase characteristics.

\textbf{Multi-agent scaling.}
The experiment is strictly two-agent. A real concurrent setting with three or
more agents operating simultaneously raises questions about notification fan-out
latency under load, correct handling of cascading writes (Agent~B receives a
directive while itself mid-write to a node that Agent~C depends on), and
registry consistency under concurrent \texttt{register\_read} and
\texttt{commit\_write} calls. These scenarios require both engineering work and
a new experimental design.

\textbf{Multi-language and \textsc{Uses\_Type} traversal.}
The current implementation and conflict classifier are Python-specific. Extending
to TypeScript or Java would require corresponding type checkers (\texttt{tsc},
\texttt{javac}) and language-specific Tree-sitter grammars. Additionally,
integrating \textsc{Uses\_Type} edge traversal as an optional BFS mode would
support type-definition-centric conflict scenarios --- cases where a type field
change rather than a function signature change is the source of the semantic
conflict --- at the cost of a higher false-positive rate that the precision
study above would help quantify.

The implementation and evaluation harness are available at
\url{https://github.com/Ibrahem-al/kovex-research}.

% =============================================================================
\section*{Acknowledgments}
% =============================================================================

\noindent\textbf{AI Use Disclosure.}
The author used Claude (Anthropic) for drafting and editing this manuscript
and for code assistance during implementation and evaluation. Perplexity AI
was used to assist in identifying relevant academic literature. All scientific
claims, formal definitions, algorithmic content, experimental design, and
results were reviewed, verified, and approved by the author. No AI system is
listed as a co-author or cited as a reference.

\bibliographystyle{abbrvnat}
\bibliography{kovex}

\end{document}
